\section{Zadanie 7}
Niech $X,Y$ będą niezależnymi zmiennymi z rozkładu $N(0,1)$. Wyznacz rozkład $X+Y$ pod warunkiem, że $X > 0$ i $Y > 0$. To znaczy, dla każdego $z$ oblicz:
\[
P(X + Y < z | X > 0, Y > 0).
\]
\textit{Wskazówka:} Wynik podaj w terminach standardowego rozkładu normalnego.

\subsection*{Rozwiązanie}

\begin{align*}
P(X + Y < z | X > 0, Y > 0) = \frac{P(X+Y<z, X>0, Y>0)}{P(X>0, Y>0)}
\end{align*}
Mianownik tego wyrażenia wynosi $ \frac{1}{2}\cdot \frac{1}{2}  =\frac{1}{4} $(z niezależności $X,Y$). \\
\\
Aby łatwiej obliczyć licznik, wprowadźmy zmienne $U = \frac{X+Y}{\sqrt{2}}, V=\frac{X-Y}{\sqrt{2}}$. Z Zadania 5(podpunkt $iii$) wiemy, że $X+Y$ oraz $X-Y$ są niezależne, więc $U$ oraz $V$ są niezależne. Oprócz tego $X+Y$ ma rozkład $N(0,2)$, a po 'przeskalowaniu' przez $\sqrt{2}$ dostajemy, że $U,V$ mają rozkład standardowy normalny.\\
Obliczmy więc licznik, po drodze korzystając z faktu, że $\phi(u)$ to pochodna $\Phi(u)$, więc $\int \phi(u)\Phi(u) du = \frac{1}{2}\Phi(u)^2$, co z kolei implikuje $\int_{0}^{t} \phi(u)\Phi(u) du = \frac{1}{2}(\Phi(t)^2 - \Phi(0)^2) = \frac{1}{2}\Phi(t)^2 - \frac{1}{8}$.
\begin{align*}
& P(X+Y<z, X>0, Y>0) = P(0<U<\frac{z}{\sqrt{2}}, -U<V<U) = \int_{0}^{\frac{z}{\sqrt{2}}} \phi(u) \int_{-u}^{u} \phi(v) dvdu = \\
&= \int_{0}^{\frac{z}{\sqrt{2}}} \phi(u) (2\Phi(u) - 1)du = 2\int_{0}^{\frac{z}{\sqrt{2}}}\phi(u)\Phi(u) du - \int_{0}^{\frac{z}{\sqrt{2}}}\phi(u)du = 
\Phi(\frac{z}{\sqrt{2}})^2 - \Phi(0)^2
- \Phi(\frac{z}{\sqrt{2}})+\frac{1}{2} =\\&= 
\Phi(\frac{z}{\sqrt{2}})^2 - \Phi(\frac{z}{\sqrt{2}}) + \frac{1}{4} = (\Phi(\frac{z}{\sqrt{2}}) - \frac{1}{2})^2
\end{align*}

Łącząc otrzymane wyniki, otrzymujemy wynik $4(\Phi(\frac{z}{\sqrt{2}}) - \frac{1}{2})^2$
\subsection*{Odpowiedź}
\[
\boxed{
4(\Phi(\frac{z}{\sqrt{2}}) - \frac{1}{2})^2 \text{ dla } z > 0\text{, wpp. } 0
}
\]
